\documentclass[11pt,a4paper]{article}
\usepackage{gtu-el-yazisi}
\begin{document}
\notbilgi{Makine Öğrenmesi}{CNN, RNN/LSTM, encoding ve TF-IDF}{Ders-9.pdf}

\sayfa{1}
\elcizimi[0.25]{sekiller/Ders-9-s1-1.png}[1]
$\underset{\text{4096'lık}}{64\times 64}$ $\rightarrow$ MLP Multilayer Perceptron

\textcolor{turuncukalem}{feature'ın yerine hassas bir ağırlıklarımız var}

\elcizimi[0.25]{sekiller/Ders-9-s1-2.png}[1]
$\leftarrow$ türev edge detection

\elcizimi[0.3]{sekiller/Ders-9-s1-3.png}[1]
bunlar konumdan bağımsız.

\elcizimi[0.25]{sekiller/Ders-9-s1-4.png}[1]

\elcizimi[0.4]{sekiller/Ders-9-s1-5.png}[1]
bu \okunamadi{W'lar} falan çalışıyor

\elcizimi[0.85]{sekiller/Ders-9-s1-6.png}[1]
featureları çıkarıyor \hfill fully connected layer

$\begin{bmatrix} -1 & 0 \\ 0 & 1 \end{bmatrix}$ \quad tek bir layer kullanılmaz

Pooling Layers $\rightarrow$ \begin{tabular}[c]{@{}l@{}}Max Pooling\\ Avg Pooling\end{tabular}

\elcizimi[0.25]{sekiller/Ders-9-s1-7.png}[1]
Feature extraction için çok kullanılır

Odaklanıp işimize yarayacağımız verilerin her şeyi zaten değil.

kernel size \quad filtre sayısı

\sayfa{2}
\elcizimi[0.7]{sekiller/Ders-9-s2-1.png}[2]

\textcircled{1} feature extraction: Öznitelik çıkarma

\textcircled{2} Özniteliklerin konumlarından bağımsız

\textcircled{3} Less parameters

eğitimi \okunamadi{ANN'ye} göre daha zor

\section{Transfer Learning}
Elimizdeki veriler ne kadar kaliteli ise o kadar kolay sınıflandırıcımız olur.

Yolo feature extraction'ı çok iyi yapıyor ve iyi parçalara ayırıyor

\elcizimi[0.4]{sekiller/Ders-9-s2-2.png}[2]

katmanlar arası isterseniz aktivation atabilirsiniz ama çok \okunamadi{kullanılmıyor}.

\section{Explainable AI}
\elcizimi[0.3]{sekiller/Ders-9-s2-3.png}[2]
time series prediction\\
Recurrent neural network (RNN)

\elcizimi[0.55]{sekiller/Ders-9-s2-4.png}[2]

Vanishing Gradient

$W_k = \gamma\,\uparrow\downarrow$ \okunamadi{} \quad Exploding Gradient

sağdan sola, soldan sağa mı sıra da önemli

Türkçedeki yüklem,

Bidirectional mimariler var

nlp\\
ses analizi

\sayfa{3}
\elcizimi[0.7]{sekiller/Ders-9-s3-1.png}[3]

Cell falan da var

\section{Encoding}
ham bir veriyi sayıya çevirmek

kedi: 0\\
köpek: 1

One hot Encoding

R: 0 \quad G: 1 \quad B: 2
\begin{gather*}
R=[1,0,0]\\
G=[0,1,0]\\
B=[0,0,1]
\end{gather*}

Ali topa at

$\rightarrow$ tokenization

Al/i/ /to/pa/ /at/.

gözetimsiz olarak entropy kullandılar.

$\text{king}-\text{man} = \text{queen}-\text{woman}$

en basit şekilde nasıl vektöre çeviririm

\sayfa{4}
\elcizimi[0.4]{sekiller/Ders-9-s4-1.png}[4]

\formulkutu{\text{TF-IDF}}
\[ \mathrm{TF}(t,d) \]
\[ \mathrm{IDF}(t) \]
Term freq (TF)

bazı çok kullanılan kelimeler var

\end{document}
